% =============================================================================
% Anonymous introduction -- three-claim framing
% =============================================================================
\section{Introduction}
\label{sec:intro}

Reinforcement learning post-training now anchors many claims about language
model reasoning and alignment~\citep{ouyang2022training,guo2025deepseekr1nature}.
PPO~\citep{schulman2017proximal}, GRPO~\citep{shao2024deepseekmath}, and
DPO~\citep{rafailov2024direct} are often compared as if the algorithm label
were the full treatment.  In practice, the outcome also depends on the base
model, instruction tuning, reward parser, group construction, sampler,
backend, LoRA configuration, checkpoint selection, and evaluation harness.

We study this problem as a systems audit rather than as an algorithm
leaderboard.  The corpus contains 79 runs across structured tool-call
emission, GSM8K, MATH-style tasks, HumanEval-style binary verification, small
and frontier-scale model families, and several training backends.  Before
presenting any results we note three deliberate scoping choices.  \emph{(i)
Algorithm fidelity.}  The training runner is a group-relative policy-gradient
variant: it retains group-normalized advantages and the critic-free structure
of GRPO, but it does not implement the PPO probability ratio and clip of
\citet{schulman2017proximal}, does not retain a frozen reference policy for
KL regularization, does one optimizer step per rollout batch, and does not
apply a completion-only token mask (a P1-A diagnostic on Qwen3.5-4B measures
$\sim$0.76 of the full-sequence loss magnitude to come from prompt tokens
rather than from the completion).  Claims in
this paper therefore refer to this runner, not to canonical GRPO.
\emph{(ii) Evaluator scope.}  The HumanEval, tool-use, and MATH rewards
measure sandboxed execution with Python built-ins removed, JSON schema
compliance, and boxed-answer formatting respectively; they are proxies, not
end-to-end capability metrics.  \emph{(iii) Probe-vs.-generalization.}  The
MATH-500 and HumanEval prompt pools in the training loop are loaded from the
\texttt{test} splits and are used as reward-environment probes, not as
held-out generalization tests.  The only genuinely held-out capability
evaluation in this paper is a separate GSM8K 200-example slice that was not
used during training.  The remaining heterogeneity in outcomes across runs is
the point: it lets us ask which failures are evidence about the policy and
which failures are evidence about the reward environment.

\paragraph{Claim 1: ZVF/GU are early triage diagnostics.}
The most reusable methodological contribution is the operational use of
Zero-Variance Fraction (ZVF) and Gradient Utilization (GU).  High ZVF with
low reward means cold-start collapse; high ZVF with high reward means
saturation; low ZVF means the current policy still samples both successes and
failures.  In a 22-run de-duplicated validation set, a fixed rule using only
the first five logged steps catches both collapsed runs with no false
positives.  This is not enough to call ZVF a calibrated predictor.  The safe
claim is narrower: ZVF/GU are cheap early-run diagnostics that convert flat
reward groups into decisions to stop, warm-start, densify reward, change task
difficulty, or increase group size.

\paragraph{Claim 2: GRPO success depends on mixed-reward groups.}
Under binary rewards, a sampled group has zero variance if every completion is
wrong or every completion is correct.  If the current policy succeeds with
success rate $p_x$ for prompt $x$ and group size is $G$, the probability of a usable mixed
group is
\[
P(\text{usable}) \approx \frac{1}{N}\sum_x \left[1 - (1-p_x)^G - p_x^G\right].
\]
The familiar $p_x > 1/G$ rule is only an expected-one-success heuristic, not a
threshold theorem.  SFT warm-up helps because it changes $p_0$; larger groups
help only when they reduce the all-wrong term more than they waste compute;
dense or partial-credit rewards help because they reduce all-zero/all-one
groups.  GRPO is therefore best understood as a reward-contrast amplifier
rather than a capability creator.

\paragraph{Claim 3: algorithm labels are incomplete treatments.}
The PPO/GRPO/DPO label is not enough to interpret an LLM post-training run.
Our Qwen PPO/GRPO comparison is artifact-sensitive: the source-of-truth ledger
records PPO at 0.225 last-10 reward, while a statistical summary records
0.350.  The Llama-3.1-8B-Instruct comparison favors PPO over GRPO.  These
observations should not be turned into a directional algorithm ranking.  They
instead show that algorithm-level claims require stack-conditioned reporting:
model family, backend, reward implementation, rollout sampler, optimizer and
KL handling, LoRA target modules, checkpoint-selection rule, and evaluation
slice must be fixed or explicitly modeled.

\paragraph{Contributions.}
This paper contributes (i) an operational ZVF/GU triage framing for standard
GRPO runs; (ii) a mixed-group probability mechanism that explains why
warm-starts, group size, reward density, and task difficulty govern whether
GRPO has signal; (iii) a source-of-truth ledger that separates online training
reward from held-out accuracy and flags artifact inconsistencies; and (iv) a
methodological warning that PPO, GRPO, and DPO comparisons in LLM
post-training are not identifiable unless the full training stack is reported.
